\section{Revalidación del umbral de decisión de Random Forest y aporte de Isolation Forest}\label{anexo:revalidacion-umbral}

Este anexo verifica si el umbral de decisión de Random Forest se mantiene estable para la generación \texttt{rf\_v11}, y mide qué aporta Isolation Forest (\texttt{if\_v10}) a la detección cuando se combina con Random Forest. Todas las cifras provienen de la evaluación sobre el conjunto de validación, no sobre el conjunto de prueba.

\subsection{Revalidación del umbral de Random Forest}

El umbral \texttt{RF\_THRESHOLD = 0.35} se fijó sobre el conjunto de validación con \texttt{rf\_v3} y se mantuvo en las generaciones posteriores sin recalibración (\texttt{docs/decision-policy.md}, §2.2.1, del repositorio logSguarDian). Para \texttt{rf\_v11} se repitió el barrido sobre el conjunto de validación de 57\,481 filas, sin cambiar el valor desplegado. El autor definió el criterio de estabilidad antes de ejecutar el barrido; no consta registro en el repositorio. El criterio exige que la variación entre los umbrales 0.30 y 0.40 sea menor que 0.5 puntos porcentuales, tanto en la tasa de bloqueo benigno como en el recall global.

Las variaciones observadas fueron de 0.0067 puntos porcentuales en la tasa de bloqueo benigno y de 0.056 puntos porcentuales en el recall global, por lo que el criterio se cumple. Esta lectura acota el resultado a estabilidad y no a optimalidad. Con un recall cercano a 99.9 \%, la validación discrimina poco entre los umbrales 0.20 y 0.40. El umbral no se recalibró, y la Configuración 2 y la Ronda 5 se mantienen con 0.35.

El Cuadro \ref{tabla:anexo-barrido-rf} presenta el barrido completo.

\begin{table}[htbp]
\centering
\footnotesize
\setlength{\tabcolsep}{3pt}
\begin{adjustbox}{max width=\linewidth}
\begin{tabular}{|r|r|r|r|r|r|r|}
\hline
\textbf{Umbral} $T$ & \textbf{Benignos bloqueados (\%)} & \textbf{Recall global (\%)} & \textbf{\texttt{sqli}} & \textbf{\texttt{xss}} & \textbf{\texttt{path\_traversal}} & \textbf{\texttt{cmdi}} \\ \hline
0.20 & 0.114 & 99.930 & 99.977 & 99.666 & 99.842 & 99.796 \\ \hline
0.25 & 0.114 & 99.930 & 99.977 & 99.666 & 99.842 & 99.796 \\ \hline
0.30 & 0.114 & 99.930 & 99.977 & 99.666 & 99.842 & 99.796 \\ \hline
\textbf{0.35} & \textbf{0.107} & \textbf{99.908} & \textbf{99.962} & \textbf{99.644} & \textbf{99.723} & \textbf{99.796} \\ \hline
0.40 & 0.107 & 99.873 & 99.938 & 99.577 & 99.644 & 99.659 \\ \hline
0.45 & 0.107 & 99.775 & 99.877 & 99.443 & 99.446 & 98.978 \\ \hline
0.50 & 0.107 & 99.634 & 99.780 & 99.354 & 99.051 & 98.093 \\ \hline
0.55 & 0.067 & 99.364 & 99.598 & 99.043 & 98.378 & 96.594 \\ \hline
0.60 & 0.067 & 99.117 & 99.417 & 98.798 & 97.785 & 95.436 \\ \hline
0.65 & 0.054 & 98.742 & 99.117 & 98.397 & 97.191 & 93.733 \\ \hline
0.70 & 0.047 & 98.296 & 98.695 & 98.063 & 96.598 & 92.643 \\ \hline
0.75 & 0.027 & 97.610 & 98.053 & 97.506 & 95.728 & 90.872 \\ \hline
0.80 & 0.020 & 96.774 & 97.218 & 96.794 & 95.055 & 89.373 \\ \hline
0.85 & 0.020 & 95.976 & 96.537 & 95.569 & 94.146 & 87.330 \\ \hline
0.90 & 0.013 & 94.858 & 95.576 & 94.166 & 92.642 & 84.128 \\ \hline
\end{tabular}
\end{adjustbox}
\caption{Barrido del umbral de decisión de Random Forest (\texttt{rf\_v11}) sobre el conjunto de validación de 57\,481 filas. Las columnas por clase son recall en porcentaje.}
\label{tabla:anexo-barrido-rf}
\par\smallskip\footnotesize Fuente: \texttt{\_evidencia/revalidacion-umbral-if/analysis\_out.txt}, sección 1. Elaboración propia.
\end{table}

\subsection{Aporte de Isolation Forest}

Con $T = 0.35$, Random Forest deja pasar 39 de los 42\,594 ataques del conjunto de validación. Isolation Forest marca 19 de esos 39 ataques: 5 de 13 de \texttt{sqli}, 10 de 16 de \texttt{xss}, 4 de 7 de \texttt{path\_traversal} y ninguno de los 3 de \texttt{cmdi} que Random Forest dejó pasar. El Cuadro \ref{tabla:anexo-aporte-if} resume estas cifras.

Se evaluó una regla hipotética que bloquea una petición si cualquiera de los dos modelos la marca. Esa regla eleva el recall global de 99.908 \% a 99.953 \%. El costo es un aumento del bloqueo benigno de 0.107 \% a 5.925 \%, es decir, 866 benignos adicionales bloqueados. El umbral de Isolation Forest (0.004206) se eligió sobre este mismo conjunto de validación, por lo que estas cifras son optimistas. La regla no se desplegó.

Estas cifras respaldan que Isolation Forest solo registre. Su aporte en detección es de 19 ataques frente a 42\,594, mientras que bloquear sus marcas agregaría 866 falsos positivos. Según el diseño del sistema, las marcas de Isolation Forest se registran con veredicto \texttt{pass\_anomaly} y no bloquean (\texttt{packages/core/src/middleware.ts}, líneas 243--244 y 333--337 del repositorio logSguarDian).

\begin{table}[htbp]
\centering
\footnotesize
\setlength{\tabcolsep}{3pt}
\begin{adjustbox}{max width=\linewidth}
\begin{tabular}{|l|r|r|}
\hline
\textbf{Concepto} & \textbf{Casos} & \textbf{Valor} \\ \hline
Ataques no bloqueados por Random Forest ($T = 0.35$) & 39 de 42\,594 & --- \\ \hline
Ataques de esos 39 marcados por Isolation Forest & 19 de 39 & --- \\ \hline
\hspace{0.5em}\texttt{sqli} & 5 de 13 & --- \\ \hline
\hspace{0.5em}\texttt{xss} & 10 de 16 & --- \\ \hline
\hspace{0.5em}\texttt{path\_traversal} & 4 de 7 & --- \\ \hline
\hspace{0.5em}\texttt{cmdi} & 0 de 3 & --- \\ \hline
Recall global, solo Random Forest & 42\,555 de 42\,594 & 99.908 \% \\ \hline
Recall global, regla ``RF o IF'' & 42\,574 de 42\,594 & 99.953 \% \\ \hline
Benignos bloqueados por Random Forest & 16 de 14\,887 & 0.107 \% \\ \hline
Benignos que Random Forest no bloquea y que Isolation Forest marca & 866 de 14\,871 & 5.82 \% \\ \hline
Benignos marcados por Isolation Forest (tasa global) & 874 de 14\,887 & 5.87 \% \\ \hline
Benignos bloqueados por la regla ``RF o IF'' & 882 de 14\,887 & 5.925 \% \\ \hline
\end{tabular}
\end{adjustbox}
\caption{Aporte de Isolation Forest sobre el conjunto de validación con $T = 0.35$ y regla hipotética ``RF o IF''.}
\label{tabla:anexo-aporte-if}
\par\smallskip\footnotesize Fuente: \texttt{\_evidencia/revalidacion-umbral-if/analysis\_out.txt}, sección 2. El numerador 874 no figura en esa salida; proviene de \texttt{\_evidencia/revalidacion-umbral-if/ua\_check\_regenerado\_out.txt} (suma de 807 y 67 marcas), y la tasa de 5.871 \% coincide con \texttt{training/models/if\_v10\_metadata.json}. Elaboración propia.
\end{table}

\subsection{Dependencia de Isolation Forest respecto del User-Agent}

De los 874 benignos que Isolation Forest marca en validación, 807 son filas sin User-Agent y 67 son filas con User-Agent. En las 14\,818 filas benignas sin User-Agent, la tasa de marca es de 5.45 \%. En cambio, entre las 69 filas benignas con User-Agent, Isolation Forest marca 67 (97.1 \%). Este grupo es pequeño, por lo que la tasa no permite una estimación precisa. Sin embargo, es compatible con una dependencia de Isolation Forest respecto de las características de User-Agent, que conviene vigilar (\texttt{\_evidencia/revalidacion-umbral-if/ua\_check\_regenerado\_out.txt}).

\subsection{Limitaciones}

Este análisis usa una sola partición de validación de 57\,481 filas. No se dispone de una segunda partición que permita medir la variación entre particiones.

El umbral de Random Forest (0.35) se eligió sobre validación con \texttt{rf\_v3}. El umbral de Isolation Forest de \texttt{if\_v10} (0.004206) se eligió sobre validación como el máximo recall entre los umbrales con recall mayor o igual a 0.50 y tasa de falsos positivos menor o igual a 0.06 (\texttt{training/retrain\_v11\_contract69.py}, líneas 107--110). Para \texttt{rf\_v11}, el gate de validación exige al menos tres de las cuatro clases de ataque con F1 mayor o igual a 0.80 (línea 60). Sus hiperparámetros, como \texttt{n\_estimators=30} y \texttt{max\_depth=25} (línea 51), están fijados en el script, que no realiza búsqueda de hiperparámetros sobre validación.

El umbral de Isolation Forest pasó de 0.00807 a 0.00421 entre las lecturas del conjunto de prueba. El umbral de 0.00807 corresponde al modelo previo al reentrenamiento, leído el 24 de septiembre de 2026. El de 0.00421 corresponde al modelo reentrenado, leído el 25 y el 26 de septiembre de 2026, y se calculó sobre validación. El historial no atribuye ese cambio al resultado de la primera lectura, pero no puede descartarse su influencia. Este hecho acota la afirmación de que el umbral se fijó solo sobre validación.

Isolation Forest no se evaluó fuera de validación en este anexo. La regla R2 de la metodología impide una nueva lectura del conjunto de prueba, por lo que la regla ``RF o IF'' no se midió sobre dicho conjunto.

Los scripts \texttt{analysis.py} y \texttt{part3.py} usan rutas absolutas de la máquina del autor, por lo que requieren ajustarlas para reproducirse.

\subsection{Procedencia}

Los modelos evaluados son \texttt{rf.onnx} e \texttt{if.onnx} del par C, descrito en la sección de pares de modelos (sección~\ref{sec:pares-modelos}). Sus sumas SHA-256 completas son las del Cuadro \ref{tabla:anexo-procedencia}. El commit del repositorio logSguarDian que produjo estos modelos es \texttt{03de50f}. El archivo \texttt{analysis\_out.txt} no registra el commit de ejecución; durante la verificación, el árbol de trabajo estaba limpio en el commit \texttt{41553a7}.

\begin{table}[htbp]
\centering
\footnotesize
\begin{adjustbox}{max width=\linewidth}
\begin{tabular}{|l|p{0.62\linewidth}|}
\hline
\textbf{Archivo} & \textbf{SHA-256} \\ \hline
\texttt{rf.onnx} (modelo, par C) & \seqsplit{25b407e649283606a231665a5a8c8ea1c2c7ad4fd0374a8945b3f119651b0870} \\ \hline
\texttt{if.onnx} (modelo, par C) & \seqsplit{5be28ee8437643b2f13708eee1ac3175348a350ea6fb8f8daeb09bd5717f797e} \\ \hline
\texttt{analysis\_out.txt} & \seqsplit{78dbf342acf01aacf112e76cd9fddfa95d3388de551a2c0eb1cd0901344835a4} \\ \hline
\texttt{part3\_out.txt} & \seqsplit{a03cb07f2e3433e7dfae4ae35866abe9828c76bffb0aa6c4afe8f049d477322e} \\ \hline
\texttt{ua\_check.txt} & \seqsplit{1027dacfbf97cc337994fbac7c989581a3fd7cd26269d1d0e5f4730081c940f7} \\ \hline
\texttt{analysis.py} (genera \texttt{analysis\_out.txt}) & \seqsplit{7c00ad61b662f899ae5d5f2f0cc73ccc5871db149d33189a6e7f16b53ea4e7d2} \\ \hline
\texttt{part3.py} (genera \texttt{part3\_out.txt}) & \seqsplit{78ace833702df2078090dad5d6d1c0bb3c83994e6cfd9e221c12b9ce5bd99d9f} \\ \hline
\texttt{ua\_check\_regenerado.py} (regenera \texttt{ua\_check.txt}) & \seqsplit{56516f5f6b63d3c3886ef4a6350807c356d4c2fc49d8ea77a0207be4adcffe5f} \\ \hline
\texttt{ua\_check\_regenerado\_out.txt} & \seqsplit{d61660015b8f763af7140987836b6b94559c025d76f2d4ff980766d5ba5fe6a6} \\ \hline
\end{tabular}
\end{adjustbox}
\caption{Sumas SHA-256 de los modelos evaluados, de los archivos de evidencia y de los scripts que los generaron o regeneraron.}
\label{tabla:anexo-procedencia}
\par\smallskip\footnotesize Fuente: hashes de los modelos, de \texttt{training/models/} en el commit \texttt{03de50f} del repositorio logSguarDian; hashes de los archivos de evidencia y de los scripts, calculados con \texttt{shasum -a 256} sobre \texttt{\_evidencia/revalidacion-umbral-if/}. El script original que generó \texttt{ua\_check.txt} no se conserva; \texttt{ua\_check\_regenerado.py} es una regeneración que reproduce las cifras de validación de ese archivo, y su salida se guarda en \texttt{ua\_check\_regenerado\_out.txt}. La línea de \texttt{ua\_check.txt} sobre el \texttt{login} no se reproduce, porque depende de \texttt{login\_vec.json}. Elaboración propia.
\end{table}

\section{Casi-duplicados entre particiones}\label{anexo:casi-duplicados}

Este anexo cuantifica la similitud textual de cada fila de validación, de prueba y de la Ronda 5 con el conjunto de entrenamiento, y la de cada fila de entrenamiento con el resto de ese conjunto. Es una auditoría posterior a la evaluación: no modifica el corpus, los modelos ni las métricas reportadas.

\subsection{Parámetros}\label{anexo:casi-parametros}

La similitud se mide sobre el texto de cada petición. Cada campo (\texttt{path}, \texttt{query} y \texttt{body}) se normaliza con decodificación \texttt{unquote\_plus}, conversión a minúsculas y colapso de espacios, y los tres campos se unen con el separador \texttt{ | }. El texto normalizado se descompone en shingles de tres caracteres, y la similitud entre dos peticiones es el índice de Jaccard exacto de sus conjuntos de shingles.

Para encontrar candidatos se usa MinHash con \texttt{num\_perm=128} y LSH con bandas de \texttt{params=(8,16)}. Cada candidato se verifica con el Jaccard exacto, de modo que todo par informado con $J \geq$ umbral es un par real; los pares que el LSH omite se tratan en las limitaciones.

Se usan los umbrales 0.8 y 0.9. Para cada fila, \textit{cualquiera} es el máximo de $J$ frente a todas las filas de entrenamiento, y \textit{misma clase} es el máximo frente a las filas con la misma etiqueta. En la Ronda 5, el texto de una petición \texttt{GET} es su \texttt{query} y el de una \texttt{POST} es su \texttt{body}, en la forma \texttt{param=payload}. Quedan fuera las filas sin texto bruto (82 de entrenamiento, 18 de validación y 19 de prueba), y el $n$ de cada cuadro lo refleja.

\subsection{Resultados}\label{anexo:casi-resultados}

Cada celda indica el número de filas con $J \geq$ umbral y su porcentaje sobre $n$. Los cuadros~\ref{tabla:anexo-casi-val}, \ref{tabla:anexo-casi-test} y \ref{tabla:anexo-casi-train} cubren las tres particiones con todas sus clases; el cuadro~\ref{tabla:anexo-casi-ronda5} cubre la Ronda 5.

\begin{table}[htbp]
\centering
\footnotesize
\setlength{\tabcolsep}{3pt}
\begin{adjustbox}{max width=\linewidth}
\begin{tabular}{|l|r|c|c|c|c|}
\hline
\textbf{Clase} & \textbf{n} & \textbf{Cualquiera $\geq$ 0.8} & \textbf{Cualquiera $\geq$ 0.9} & \textbf{Misma clase $\geq$ 0.8} & \textbf{Misma clase $\geq$ 0.9} \\ \hline
\texttt{benign} & 14\,887 & 4\,491 (30.17\,\%) & 3\,241 (21.77\,\%) & 4\,490 (30.16\,\%) & 3\,241 (21.77\,\%) \\ \hline
\texttt{sqli} & 34\,103 & 25\,356 (74.35\,\%) & 18\,840 (55.24\,\%) & 25\,308 (74.21\,\%) & 18\,804 (55.14\,\%) \\ \hline
\texttt{xss} & 4\,484 & 2\,076 (46.30\,\%) & 1\,115 (24.87\,\%) & 2\,054 (45.81\,\%) & 1\,112 (24.80\,\%) \\ \hline
\texttt{path\_traversal} & 2\,526 & 1\,751 (69.32\,\%) & 1\,287 (50.95\,\%) & 1\,704 (67.46\,\%) & 1\,253 (49.60\,\%) \\ \hline
\texttt{cmdi} & 1\,463 & 580 (39.64\,\%) & 323 (22.08\,\%) & 542 (37.05\,\%) & 302 (20.64\,\%) \\ \hline
\textbf{Total} & 57\,463 & 34\,254 (59.61\,\%) & 24\,806 (43.17\,\%) & 34\,098 (59.34\,\%) & 24\,712 (43.01\,\%) \\ \hline
\end{tabular}
\end{adjustbox}
\caption{Casi-duplicados de la partición de validación frente al entrenamiento (57\,463 filas con texto).}
\label{tabla:anexo-casi-val}
\par\smallskip\footnotesize Fuente: \texttt{neardup\_v2\_results.json}, sección \texttt{val}; \texttt{tables.md}. Elaboración propia.
\end{table}


\begin{table}[htbp]
\centering
\footnotesize
\setlength{\tabcolsep}{3pt}
\begin{adjustbox}{max width=\linewidth}
\begin{tabular}{|l|r|c|c|c|c|}
\hline
\textbf{Clase} & \textbf{n} & \textbf{Cualquiera $\geq$ 0.8} & \textbf{Cualquiera $\geq$ 0.9} & \textbf{Misma clase $\geq$ 0.8} & \textbf{Misma clase $\geq$ 0.9} \\ \hline
\texttt{benign} & 14\,886 & 4\,404 (29.58\,\%) & 3\,176 (21.34\,\%) & 4\,403 (29.58\,\%) & 3\,175 (21.33\,\%) \\ \hline
\texttt{sqli} & 34\,103 & 25\,569 (74.98\,\%) & 18\,841 (55.25\,\%) & 25\,511 (74.81\,\%) & 18\,809 (55.15\,\%) \\ \hline
\texttt{xss} & 4\,484 & 2\,061 (45.96\,\%) & 1\,090 (24.31\,\%) & 2\,038 (45.45\,\%) & 1\,087 (24.24\,\%) \\ \hline
\texttt{path\_traversal} & 2\,526 & 1\,788 (70.78\,\%) & 1\,314 (52.02\,\%) & 1\,749 (69.24\,\%) & 1\,290 (51.07\,\%) \\ \hline
\texttt{cmdi} & 1\,463 & 557 (38.07\,\%) & 326 (22.28\,\%) & 523 (35.75\,\%) & 299 (20.44\,\%) \\ \hline
\textbf{Total} & 57\,462 & 34\,379 (59.83\,\%) & 24\,747 (43.07\,\%) & 34\,224 (59.56\,\%) & 24\,660 (42.92\,\%) \\ \hline
\end{tabular}
\end{adjustbox}
\caption{Casi-duplicados de la partición de prueba frente al entrenamiento (57\,462 filas con texto). Solo se comparó el texto; el conjunto de prueba no se puntuó.}
\label{tabla:anexo-casi-test}
\par\smallskip\footnotesize Fuente: \texttt{neardup\_v2\_results.json}, sección \texttt{test}; \texttt{tables.md}. Elaboración propia.
\end{table}


\begin{table}[htbp]
\centering
\footnotesize
\setlength{\tabcolsep}{3pt}
\begin{adjustbox}{max width=\linewidth}
\begin{tabular}{|l|r|c|c|c|c|}
\hline
\textbf{Clase} & \textbf{n} & \textbf{Cualquiera $\geq$ 0.8} & \textbf{Cualquiera $\geq$ 0.9} & \textbf{Misma clase $\geq$ 0.8} & \textbf{Misma clase $\geq$ 0.9} \\ \hline
\texttt{benign} & 69\,472 & 20\,765 (29.89\,\%) & 14\,824 (21.34\,\%) & 20\,763 (29.89\,\%) & 14\,822 (21.34\,\%) \\ \hline
\texttt{sqli} & 159\,138 & 118\,649 (74.56\,\%) & 88\,043 (55.32\,\%) & 118\,375 (74.39\,\%) & 87\,878 (55.22\,\%) \\ \hline
\texttt{xss} & 20\,929 & 9\,534 (45.55\,\%) & 5\,217 (24.93\,\%) & 9\,475 (45.27\,\%) & 5\,201 (24.85\,\%) \\ \hline
\texttt{path\_traversal} & 11\,787 & 8\,257 (70.05\,\%) & 6\,035 (51.20\,\%) & 8\,044 (68.24\,\%) & 5\,930 (50.31\,\%) \\ \hline
\texttt{cmdi} & 6\,832 & 2\,697 (39.48\,\%) & 1\,508 (22.07\,\%) & 2\,508 (36.71\,\%) & 1\,393 (20.39\,\%) \\ \hline
\textbf{Total} & 268\,158 & 159\,902 (59.63\,\%) & 115\,627 (43.12\,\%) & 159\,165 (59.35\,\%) & 115\,224 (42.97\,\%) \\ \hline
\end{tabular}
\end{adjustbox}
\caption{Casi-duplicados del entrenamiento frente a sí mismo, excluida la propia fila (268\,158 filas con texto).}
\label{tabla:anexo-casi-train}
\par\smallskip\footnotesize Fuente: \texttt{neardup\_v2\_results.json}, sección \texttt{train\_self}; \texttt{tables.md}. Elaboración propia.
\end{table}


\begin{table}[htbp]
\centering
\footnotesize
\setlength{\tabcolsep}{3pt}
\begin{adjustbox}{max width=\linewidth}
\begin{tabular}{|l|r|c|c|c|c|}
\hline
\textbf{Clase} & \textbf{n} & \textbf{Cualquiera $\geq$ 0.8} & \textbf{Cualquiera $\geq$ 0.9} & \textbf{Misma clase $\geq$ 0.8} & \textbf{Misma clase $\geq$ 0.9} \\ \hline
\texttt{sqli} & 200 & 0 (0.00\,\%) & 0 (0.00\,\%) & 0 (0.00\,\%) & 0 (0.00\,\%) \\ \hline
\texttt{path\_traversal} & 200 & 0 (0.00\,\%) & 0 (0.00\,\%) & 0 (0.00\,\%) & 0 (0.00\,\%) \\ \hline
\texttt{cmdi} & 200 & 21 (10.50\,\%) & 7 (3.50\,\%) & 21 (10.50\,\%) & 7 (3.50\,\%) \\ \hline
\textbf{Total} & 600 & 21 (3.50\,\%) & 7 (1.17\,\%) & 21 (3.50\,\%) & 7 (1.17\,\%) \\ \hline
\end{tabular}
\end{adjustbox}
\caption{Casi-duplicados de la Ronda 5 (200 peticiones por clase) frente al entrenamiento.}
\label{tabla:anexo-casi-ronda5}
\par\smallskip\footnotesize Fuente: \texttt{neardup\_v2\_results.json}, sección \texttt{round5}; \texttt{tables.md}. Elaboración propia.
\end{table}

Con umbral 0.8, el 59.61\,\% de las filas de validación, el 59.83\,\% de las de prueba y el 59.63\,\% de las de entrenamiento tienen al menos un vecino de cualquier clase con $J \geq 0.8$. Con la misma clase, esas cifras son 59.34\,\%, 59.56\,\% y 59.35\,\%. En la Ronda 5, 21 de las 200 peticiones de \texttt{cmdi} tienen un casi-duplicado en entrenamiento con $J \geq 0.8$; ninguna de \texttt{sqli} ni de \texttt{path\_traversal} lo tiene.

\subsection{Desempeño en validación según haya o no vecino en entrenamiento}\label{anexo:casi-validacion-vecino}

El cuadro separa las filas de validación según tengan o no un vecino en entrenamiento con $J \geq 0.8$. Reporta la tasa de bloqueo de benignos (falsos positivos) y el recall por clase, con su intervalo de Wilson al 95\,\%. No se reporta F1.

\begin{table}[htbp]
\centering
\footnotesize
\begin{adjustbox}{max width=\linewidth}
\begin{tabular}{|l|l|r|c|c|}
\hline
\textbf{Grupo} & \textbf{Clase} & \textbf{Bloqueados/n} & \textbf{Tasa} & \textbf{IC 95\,\% (Wilson)} \\ \hline
Con vecino & Benignos (tasa de bloqueo, FP) & 1/4\,491 & 0.0002 & [0.0000, 0.0013] \\ \hline
Con vecino & \texttt{cmdi} (recall) & 578/580 & 0.9966 & [0.9875, 0.9991] \\ \hline
Con vecino & \texttt{path\_traversal} (recall) & 1\,750/1\,751 & 0.9994 & [0.9968, 0.9999] \\ \hline
Con vecino & \texttt{sqli} (recall) & 25\,354/25\,356 & 0.9999 & [0.9997, 1.0000] \\ \hline
Con vecino & \texttt{xss} (recall) & 2\,074/2\,076 & 0.9990 & [0.9965, 0.9997] \\ \hline
Sin vecino & Benignos (tasa de bloqueo, FP) & 15/10\,396 & 0.0014 & [0.0009, 0.0024] \\ \hline
Sin vecino & \texttt{cmdi} (recall) & 882/883 & 0.9989 & [0.9936, 0.9998] \\ \hline
Sin vecino & \texttt{path\_traversal} (recall) & 770/775 & 0.9935 & [0.9850, 0.9972] \\ \hline
Sin vecino & \texttt{sqli} (recall) & 8\,736/8\,747 & 0.9987 & [0.9977, 0.9993] \\ \hline
Sin vecino & \texttt{xss} (recall) & 2\,394/2\,408 & 0.9942 & [0.9903, 0.9965] \\ \hline
\end{tabular}
\end{adjustbox}
\caption{Bloqueos y recall en validación, separados por la existencia de un vecino en entrenamiento con $J \geq 0.8$. Modelo par C (\texttt{rf.onnx} e \texttt{if.onnx}; sin F1).}
\label{tabla:anexo-casi-vecino}
\par\smallskip\footnotesize Fuente: \texttt{block\_gap\_counts.json} y \texttt{block\_gap\_counts.md} en \texttt{\_evidencia/casi-duplicados-v3/}. Elaboración propia.
\end{table}

Lectura acotada. Las diferencias de recall entre ambos grupos son de 0.6 puntos porcentuales o menos en todas las clases; la mayor es de 0.59 en \texttt{path\_traversal}. De los 16 falsos positivos de validación, 15 están en filas sin vecino. Con tan pocos casos, la tasa de ese grupo no se estima con precisión. El grupo ``sin vecino'' contiene además filas que sí tienen un vecino real que el LSH no encontró, cerca del 11\,\% de los pares con $J \geq 0.8$ según la medición de recall de este anexo. Por eso la brecha real entre grupos puede ser mayor que la observada. Estos datos no establecen una relación causal.

\subsection{Cobertura de \texttt{cmdi} en la Ronda 5 sin los payloads con vecino}\label{anexo:casi-cobertura-cmdi}

La cobertura de \texttt{cmdi} en la Ronda 5 se recalcula sin los payloads que tienen un vecino en entrenamiento. El denominador es de 200 payloads, y los timeouts, errores y respuestas HTTP 5xx cuentan como no bloqueados. El par A usa \texttt{results\_config2.json} del commit \texttt{a3aada0}, y el par C usa \texttt{results\_config2-develop-tarball-20261009.json}, añadido en el commit \texttt{4d928b5}.

\begin{table}[htbp]
\centering
\footnotesize
\begin{adjustbox}{max width=\linewidth}
\begin{tabular}{|l|l|r|c|c|}
\hline
\textbf{Par} & \textbf{Subconjunto} & \textbf{Bloqueados/n} & \textbf{Cobertura} & \textbf{IC 95\,\% (Wilson)} \\ \hline
C & Todos los payloads & 192/200 & 0.9600 & [0.923, 0.980] \\ \hline
C & Sin $J \geq 0.8$ (excluye 21) & 171/179 & 0.9553 & [0.914, 0.977] \\ \hline
C & Sin $J \geq 0.9$ (excluye 7) & 185/193 & 0.9585 & [0.920, 0.979] \\ \hline
A & Todos los payloads & 183/200 & 0.9150 & [0.868, 0.946] \\ \hline
A & Sin $J \geq 0.8$ (excluye 21) & 162/179 & 0.9050 & [0.853, 0.940] \\ \hline
A & Sin $J \geq 0.9$ (excluye 7) & 176/193 & 0.9119 & [0.863, 0.944] \\ \hline
\end{tabular}
\end{adjustbox}
\caption{Cobertura de \texttt{cmdi} en la Ronda 5, con y sin los payloads que tienen un vecino en entrenamiento.}
\label{tabla:anexo-casi-cmdi-ronda5}
\par\smallskip\footnotesize Fuente: \texttt{cmdi\_coverage\_v2.json} y \texttt{cmdi\_coverage\_v2.md} en \texttt{\_evidencia/casi-duplicados-v3/}, sobre los resultados de \texttt{results\_config2.json} (commit \texttt{a3aada0}) y \texttt{results\_config2-develop-tarball-20261009.json} (commit \texttt{4d928b5}). Elaboración propia.
\end{table}

Lectura acotada. Los 21 payloads de \texttt{cmdi} con vecino a $J \geq 0.8$ están todos bloqueados en ambos pares. Por eso, al excluirlos, la cobertura baja: de 96.0\,\% a 95.5\,\% en el par C, y de 91.5\,\% a 90.5\,\% en el par A. Los intervalos de los subconjuntos son amplios por su tamaño (179 y 193 payloads).

\subsection{Limitaciones}\label{anexo:casi-limitaciones}

Las cifras son cotas inferiores. El LSH puede omitir candidatos. Se midió su recall sobre una muestra aleatoria de 2\,000 filas de validación (semilla 42): de 1\,351 filas con un par exacto de $J \geq 0.8$, el LSH encontró 1\,206 (89.3\,\%); con $J \geq 0.9$, encontró 889 de 911 (97.6\,\%). Esa medición cubre solo validación; no se midió el recall sobre prueba ni sobre entrenamiento.

Otras limitaciones:
\begin{itemize}
\item Se usa una sola normalización del texto y un solo tamaño de shingle (3 caracteres). Otra normalización cambiaría los valores de $J$.
\item El muestreo de recall se hizo con un solo corte aleatorio.
\item Solo se comparó el texto de prueba con el de entrenamiento. No se ejecutó inferencia del modelo sobre las filas casi duplicadas.
\end{itemize}

\subsection{Lectura}\label{anexo:casi-lectura}

Las particiones comparten la estructura de casi-duplicados del conjunto de entrenamiento; las métricas de prueba miden desempeño en esa distribución. Este anexo describe la estructura y no establece su causa.

\subsection{Procedencia}\label{anexo:casi-procedencia}

\begin{table}[htbp]
\centering
\footnotesize
\begin{adjustbox}{max width=\linewidth}
\begin{tabular}{|l|p{0.62\linewidth}|}
\hline
\textbf{Archivo} & \textbf{SHA-256} \\ \hline
\texttt{neardup\_v2.py} (script) & \seqsplit{d0879391e85e9a3d2972cab462d6207c8a62584dc40953ea30c9c2e8a17e52b7} \\ \hline
\texttt{neardup\_v2\_results.json} & \seqsplit{f9d56b77de0e3d8788a8de3e9d179fbd7bb2324801c72e4f08ef3e40b355a8c0} \\ \hline
\texttt{tables.md} & \seqsplit{d3620c8a155f7b9ea3aa6a75762c794509208b6475e4c177f049f4679190d8f5} \\ \hline
\texttt{run.log} & \seqsplit{e544d48f603320008941781cdef2b55a5530e6d4fc4a5f84dfa69fadc826ccb8} \\ \hline
\texttt{train.parquet} (entrada) & \seqsplit{979ad08b629b94e974fd80e1fa9ef8f4f146fc800151ca6cb8515898045b5592} \\ \hline
\texttt{val.parquet} (entrada) & \seqsplit{42b6a546e672d9961e45cf75895520590267bfdeeaaa8c4df95ab6e032970ad1} \\ \hline
\texttt{test.parquet} (entrada) & \seqsplit{3e93b7763732cbd5e23738fba5514d64b826d37635ccdf7f6ee941bce3cfe608} \\ \hline
\texttt{unified.jsonl} (entrada) & \seqsplit{3e74be1b64daf292b44adbd1f5472858967a0079af0cdb8748771d733a0951b2} \\ \hline
\texttt{large\_corpus\_requests.json} (entrada) & \seqsplit{32a171803259cfc4cde87cfb3cf4e740a4eef26d13b562d885ea8abdf7348b1a} \\ \hline
\texttt{cmdi\_coverage\_v2.py} (\texttt{casi-duplicados-v3}) & \seqsplit{038476a29e4b640036b3c94b9f1be382f7ed8594d46d31ae028d9ba19f260c60} \\ \hline
\texttt{block\_gap\_counts.py} (\texttt{casi-duplicados-v3}) & \seqsplit{bee2b350cad08a02af21de2be562106901069c654c5751cbf3185b5a23de9844} \\ \hline
\texttt{cmdi\_coverage\_v2.json} (\texttt{casi-duplicados-v3}) & \seqsplit{27c18067f10a4263d3037520da1ee1f1bd817d11b25220e6d9509c6c30c8c9bc} \\ \hline
\texttt{cmdi\_coverage\_v2.md} (\texttt{casi-duplicados-v3}) & \seqsplit{0c4709e9d175dbd2111bc462d7ba8be122aca45d4273fd5d5de0418737856db8} \\ \hline
\texttt{block\_gap\_counts.json} (\texttt{casi-duplicados-v3}) & \seqsplit{64e1f288c0354247e1fb6844b16409bda3ada23cae2b7bb61e5b023c738142a8} \\ \hline
\texttt{block\_gap\_counts.md} (\texttt{casi-duplicados-v3}) & \seqsplit{2f07a2393b45ada764a032fc54fe4655e6a3809921db6c5f2ede81cfb0e01017} \\ \hline
\texttt{provenance.md} (\texttt{casi-duplicados-v3}) & \seqsplit{ab60e7da6227b5107a40bf950ddb4e0e77c28975120c7c6bfb7683090bad4c76} \\ \hline
\end{tabular}
\end{adjustbox}
\caption{Sumas SHA-256 del script, de los resultados de la auditoría de casi-duplicados y de sus entradas.}
\label{tabla:anexo-casi-procedencia}
\par\smallskip\footnotesize Fuente: hashes del script, de \texttt{neardup\_v2\_results.json}, \texttt{tables.md} y \texttt{run.log} en \texttt{\_evidencia/casi-duplicados/}, y de las entradas según el registro de la ejecución del 11 de octubre de 2026 en \texttt{neardup\_v2\_results.json}. El hash de \texttt{large\_corpus\_requests.json} coincide con el archivo del commit \texttt{58033a0} del repositorio logSguarDian-vulnerable-project. La ejecución se hizo sobre el árbol de trabajo del repositorio logSguarDian en el commit \texttt{41553a7}, que estaba limpio. Ese commit es ancestro de \texttt{origin/develop} (verificado tras \texttt{git fetch origin}) y figura en la rama remota \texttt{origin/audit/config2-breakdown}. Los archivos \texttt{train}, \texttt{val}, \texttt{test} y \texttt{unified.jsonl} no están versionados en git en ese commit; se identifican por su SHA-256. Elaboración propia.
\end{table}

\section{Generalización entre fuentes (exclusión por fuente con la configuración de rf\_v11)}\label{anexo:loso}

Este anexo mide cuánto depende el desempeño del modelo de fuentes concretas del corpus. Para cada fuente $S$, se entrena un Random Forest con la configuración de \texttt{rf\_v11} sobre las filas de entrenamiento que no pertenecen a $S$, y se evalúa sobre las filas de $S$. El modelo de cada fold no es el \texttt{rf.onnx} evaluado en el capítulo de Resultados.

\subsection{Protocolo}\label{anexo:loso-protocolo}

El Random Forest usa los valores de \texttt{retrain\_v11\_contract69.py} (líneas 50--51 en \texttt{origin/develop}):

\begin{itemize}
\item \texttt{n\_estimators=30};
\item \texttt{max\_depth=25};
\item \texttt{class\_weight=balanced\_subsample};
\item \texttt{random\_state=42}.
\end{itemize} Las características son las 69 del contrato de producción, leídas de los archivos \texttt{parquet} ya extraídos (regla R1). El archivo \texttt{test.parquet} no se abrió (regla R2).

Se auditan las fuentes con al menos 30 filas en train + val (\texttt{MIN\_ROWS = 30}, \texttt{loso\_rf\_v11.py}). De las 17 fuentes del corpus, 15 cumplen ese mínimo y son las que aparecen en los cuadros. Las dos restantes son \texttt{synthetic\_evasion} (20 filas en train + val) y \texttt{patt\_cmdi\_curated} (13 filas), según el recuento de \texttt{source\_counts.txt} (cuadro de procedencia). Ambas tienen menos de 30 filas, por lo que el script no las audita. El comentario de \texttt{training/label\_map.yaml} indica 21 payloads para \texttt{patt\_cmdi\_curated} y el recuento de train y val da 13; no consta el motivo.

\subsection{Resultados por fuente}\label{anexo:loso-resultados}

El cuadro~\ref{tabla:loso-fuentes} reporta, para las 15 fuentes, el bloqueo de sus ataques y la tasa de falsos positivos sobre sus benignos, con todas sus filas en train + val. Las fuentes con $n < 200$ se marcan y sus intervalos son amplios.

\begin{table}[htbp]
\centering
\footnotesize
\begin{adjustbox}{max width=\linewidth}
\begin{tabular}{|l|r|c|r|c|r|c|}
\hline
\textbf{Fuente} & \textbf{n} & \textbf{n<200} & \textbf{Ataques} & \textbf{Bloqueo de ataques [IC 95\,\%]} & \textbf{Benignos} & \textbf{FP [IC 95\,\%]} \\ \hline
\texttt{capec} & 181\,486 &  & 181\,486 & 0.960 [0.959, 0.961] & 0 & — \\ \hline
\texttt{modsec\_learn} & 84\,698 &  & 25\,962 & 0.996 [0.996, 0.997] & 58\,736 & 0.753 [0.750, 0.757] \\ \hline
\texttt{payload\_full} & 25\,262 &  & 9\,959 & 0.740 [0.731, 0.748] & 15\,303 & 0.003 [0.002, 0.003] \\ \hline
\texttt{payloads\_csv} & 15\,650 &  & 10\,146 & 0.908 [0.902, 0.913] & 5\,504 & 0.266 [0.255, 0.278] \\ \hline
\texttt{xss\_dataset} & 9\,226 &  & 6\,224 & 0.995 [0.993, 0.996] & 3\,002 & 0.949 [0.941, 0.957] \\ \hline
\texttt{owasp\_logs} & 3\,280 &  & 3\,280 & 0.917 [0.907, 0.926] & 0 & — \\ \hline
\texttt{seclists\_cmdi} & 2\,087 &  & 2\,087 & 0.288 [0.269, 0.308] & 0 & — \\ \hline
\texttt{command\_injection} & 1\,745 &  & 401 & 0.509 [0.460, 0.557] & 1\,344 & 0.000 [0.000, 0.003] \\ \hline
\texttt{synthetic\_cmdi\_minimal\_context} & 709 &  & 709 & 0.635 [0.599, 0.669] & 0 & — \\ \hline
\texttt{synthetic\_cmdi\_compound} & 425 &  & 425 & 1.000 [0.991, 1.000] & 0 & — \\ \hline
\texttt{synthetic\_nav} & 357 &  & 0 & — & 357 & 0.739 [0.692, 0.782] \\ \hline
\texttt{synthetic\_cmdi\_windows} & 340 &  & 340 & 1.000 [0.989, 1.000] & 0 & — \\ \hline
\texttt{synthetic\_multifield} & 210 &  & 210 & 0.905 [0.857, 0.938] & 0 & — \\ \hline
\texttt{synthetic\_nav\_ecommerce} & 113 & sí & 0 & — & 113 & 0.788 [0.703, 0.853] \\ \hline
\texttt{mlops\_telemetry\_curated} & 100 & sí & 100 & 0.970 [0.915, 0.990] & 0 & — \\ \hline
\end{tabular}
\end{adjustbox}
\caption{Exclusión por fuente: bloqueo de ataques y falsos positivos benignos sobre todas las filas de cada fuente (train + val), con intervalo de Wilson al 95\,\%.}
\label{tabla:loso-fuentes}
\par\smallskip\footnotesize Fuente: \texttt{results.json} y \texttt{tables\_v5.md} (\texttt{\_evidencia/loso-rf-v11/}, corrida \texttt{run2}). Elaboración propia.
\end{table}

El cuadro~\ref{tabla:loso-control} compara, para las ocho fuentes grandes y sobre sus filas de validación, el modelo de la exclusión (c) con el \texttt{rf.onnx} (a) y con un control (b) que quita de entrenamiento un número aleatorio de filas igual al de la fuente, con cinco semillas.

\begin{table}[htbp]
\centering
\footnotesize
\begin{adjustbox}{max width=\linewidth}
\begin{tabular}{|l|r|c|c|c|c|c|c|}
\hline
\textbf{Fuente} & \textbf{n val} & \textbf{(a) bloqueo} & \textbf{(a) FP} & \textbf{(b) bloqueo} & \textbf{(b) FP} & \textbf{(c) bloqueo} & \textbf{(c) FP} \\ \hline
\texttt{capec} & 32\,028 & 1.000 [1.000, 1.000] & — & 1.000 [1.000, 1.000] & — & 0.960 [0.958, 0.962] & — \\ \hline
\texttt{modsec\_learn} & 14\,947 & 0.998 [0.997, 0.999] & 0.001 [0.000, 0.001] & 0.998 [0.998, 0.998] & 0.001 [0.001, 0.001] & 0.997 [0.995, 0.998] & 0.760 [0.752, 0.768] \\ \hline
\texttt{payload\_full} & 4\,459 & 0.998 [0.994, 0.999] & 0.000 [0.000, 0.002] & 0.997 [0.997, 0.998] & 0.000 [0.000, 0.000] & 0.737 [0.716, 0.757] & 0.001 [0.000, 0.003] \\ \hline
\texttt{payloads\_csv} & 2\,761 & 0.998 [0.994, 0.999] & 0.004 [0.002, 0.011] & 0.997 [0.996, 0.998] & 0.005 [0.004, 0.006] & 0.906 [0.892, 0.919] & 0.270 [0.243, 0.299] \\ \hline
\texttt{xss\_dataset} & 1\,628 & 0.990 [0.982, 0.994] & 0.002 [0.000, 0.011] & 0.991 [0.991, 0.993] & 0.000 [0.000, 0.000] & 0.993 [0.986, 0.996] & 0.945 [0.923, 0.962] \\ \hline
\texttt{owasp\_logs} & 578 & 0.990 [0.978, 0.995] & — & 0.994 [0.993, 0.995] & — & 0.915 [0.890, 0.935] & — \\ \hline
\texttt{seclists\_cmdi} & 368 & 1.000 [0.990, 1.000] & — & 1.000 [1.000, 1.000] & — & 0.299 [0.254, 0.348] & — \\ \hline
\texttt{command\_injection} & 308 & 0.958 [0.883, 0.986] & 0.000 [0.000, 0.016] & 0.952 [0.944, 0.972] & 0.000 [0.000, 0.000] & 0.563 [0.448, 0.673] & 0.000 [0.000, 0.016] \\ \hline
\end{tabular}
\end{adjustbox}
\caption{Validación de cada fuente grande: (a) \texttt{rf.onnx} par C; (b) control de mismo tamaño, media [mín, máx] de 5 semillas; (c) exclusión de la fuente. Los corchetes de (a) y (c) son intervalos de Wilson al 95\,\%.}
\label{tabla:loso-control}
\par\smallskip\footnotesize Fuente: \texttt{tables\_v5.md}, sección 2 (\texttt{\_evidencia/loso-rf-v11/}). El control se describe en \texttt{control\_wilson\_audit.py}. Elaboración propia.
\end{table}

\subsection{Composición de payload\_full}\label{anexo:loso-payload-full}

El cuadro~\ref{tabla:loso-payload-full} muestra la clase que el modelo de exclusión de \texttt{payload\_full} asigna a sus filas \texttt{sqli}. El 74.2\,\% se asigna a \texttt{cmdi}, de las cuales el 99.4\,\% se bloquea. El 25.8\,\% se asigna a \texttt{benign}, y ninguna se bloquea. Este cuadro describe la distribución de predicciones y no atribuye causa a ella.

\begin{table}[htbp]
\centering
\footnotesize
\begin{adjustbox}{max width=\linewidth}
\begin{tabular}{|l|r|c|c|}
\hline
\textbf{Clase predicha (argmax)} & \textbf{Filas} & \textbf{\% de sqli [IC]} & \textbf{Bloqueadas / predichas [IC]} \\ \hline
\texttt{benign} & 2\,378 & 0.258 [0.249, 0.267] & 0.000 [0.000, 0.002] \\ \hline
\texttt{cmdi} & 6\,846 & 0.742 [0.733, 0.751] & 0.994 [0.992, 0.996] \\ \hline
\end{tabular}
\end{adjustbox}
\caption{\texttt{payload\_full}: filas \texttt{sqli} de train + val ($n = 9\,224$) según la clase que predice el modelo de exclusión.}
\label{tabla:loso-payload-full}
\par\smallskip\footnotesize Fuente: \texttt{tables\_v5.md}, sección 3. Elaboración propia.
\end{table}

\subsection{Lectura acotada}\label{anexo:loso-lectura}

Sin causas, con las cifras de los cuadros~\ref{tabla:loso-control} y~\ref{tabla:loso-fuentes}. El bloqueo de ataques de (c) frente a (a), sobre las filas de cada fuente en validación, cae de forma distinta según la fuente. Caída grande, de más de 25 puntos porcentuales: \texttt{payload\_full} (0.998 a 0.737), \texttt{seclists\_cmdi} (1.000 a 0.299) y \texttt{command\_injection} (0.958 a 0.563). Caída moderada, de 7 a 10 puntos: \texttt{owasp\_logs} (0.990 a 0.915, 7.5 puntos) y \texttt{payloads\_csv} (0.998 a 0.906, 9.2 puntos). Caída leve, de 4 puntos: \texttt{capec} (1.000 a 0.960). Sin cambio: \texttt{modsec\_learn} (0.998 a 0.997) y \texttt{xss\_dataset} (0.990 a 0.993). Los falsos positivos benignos suben al excluir \texttt{modsec\_learn} (0.001 a 0.760), \texttt{payloads\_csv} (0.004 a 0.270) y \texttt{xss\_dataset} (0.002 a 0.945). El control de mismo tamaño (b) no reproduce estas caídas: su bloqueo queda entre 0.952 y 1.000 en las ocho fuentes.

\subsection{Limitaciones}\label{anexo:loso-limitaciones}

\begin{itemize}
\item El modelo de cada fold no es el \texttt{rf.onnx} evaluado. Los valores de (c) describen modelos de exclusión, no el sistema desplegado.
\item No se midieron casi-duplicados entre fuentes. Una fuente puede tener vecinos casi idénticos en otra, lo que afecta la lectura de (c).
\item \texttt{capec} concentra la mayor parte de \texttt{sqli} y comparte \texttt{path\_traversal} con \texttt{owasp\_logs}. Excluirla retira una parte grande de la señal de esas clases, según la composición documentada en \texttt{docs/results.md} (sección F3.6).
\item Hay fuentes con $n < 200$ (\texttt{mlops\_telemetry\_curated}, \texttt{synthetic\_nav\_ecommerce}), cuyos intervalos son amplios.
\item Se usa una sola configuración de hiperparámetros. Las cinco semillas se aplican solo al control (b).
\item Solo se evalúa Random Forest. No se reporta macro F1 ni se compara con \texttt{rf\_v7}.
\end{itemize}

\subsection{Procedencia}\label{anexo:loso-procedencia}

\begin{table}[htbp]
\centering
\footnotesize
\begin{adjustbox}{max width=\linewidth}
\begin{tabular}{|l|p{0.62\linewidth}|}
\hline
\textbf{Archivo} & \textbf{SHA-256} \\ \hline
\texttt{loso\_rf\_v11.py} & \seqsplit{e9e7484715471bf35c83aabcffd5e8cc7f74a123ed34117928ca9c2bd23c4db6} \\ \hline
\texttt{control\_wilson\_audit.py} & \seqsplit{1dbbe7db7628a11c8cfd6de79ca4f7338df06d6994a68ab0006b2f886e9a9a80} \\ \hline
\texttt{make\_tables\_v5.py} & \seqsplit{48ac94e93436394d936d89d374e4a916892035332509d8660789c25fe57c5c29} \\ \hline
\texttt{results.json (salida, \texttt{run2})} & \seqsplit{192811fb4900a030da8ab602f75cb9bdcc14c782469a27a4fc145370fe6be8f4} \\ \hline
\texttt{tables\_v5.md (salida)} & \seqsplit{574f4b2b73940f74ef705d7790fa6961b2a2735dce8ad5bd46441278cef1bd4e} \\ \hline
\texttt{sha256\_outputs\_v5.txt} & \seqsplit{03bd017abcf9236253561e4ddd1c4e8842b440fa7b947e8aa96e41d7f56ebecb} \\ \hline
\texttt{sha256\_inputs.txt} & \seqsplit{53dc7a69695c36885bf7846d4843162ab5a2707f75a3179bdc924f5b57e2a2be} \\ \hline
\texttt{train.parquet (entrada)} & \seqsplit{979ad08b629b94e974fd80e1fa9ef8f4f146fc800151ca6cb8515898045b5592} \\ \hline
\texttt{val.parquet (entrada)} & \seqsplit{42b6a546e672d9961e45cf75895520590267bfdeeaaa8c4df95ab6e032970ad1} \\ \hline
\texttt{rf.onnx (entrada, par C)} & \seqsplit{25b407e649283606a231665a5a8c8ea1c2c7ad4fd0374a8945b3f119651b0870} \\ \hline
\texttt{loso\_results.json (referencia rf\_v7)} & \seqsplit{27672b0c7e5f78fa0fee6a602fc6dd0e9a8a9f8d0a0dc90d8ae00c6acab63646} \\ \hline
\texttt{source\_counts.txt} (recuento por fuente, \texttt{python3 count\_sources.py}) & \seqsplit{b4f419e2ff83e1df15fc73a642b6e7b9cf1c780624cb9f9f7aed593fd9843289} \\ \hline
\texttt{count\_sources.py} (script del recuento) & \seqsplit{83789ced7c97921aef09452278e4f80098e7477bd7d67ece227c5000ac8ea29b} \\ \hline
\end{tabular}
\end{adjustbox}
\caption{Sumas SHA-256 de los scripts, de las salidas y de las entradas de la auditoría por exclusión de fuente.}
\label{tabla:loso-procedencia}
\par\smallskip\footnotesize Fuente: hashes de \texttt{\_evidencia/loso-rf-v11/} (\texttt{sha256\_inputs.txt}, \texttt{sha256\_outputs\_v5.txt}) y de los archivos de la ejecución. Código de referencia: commit \texttt{dc380e3} de logSguarDian (\texttt{origin/develop}); la corrida repetida reproduce los valores de la corrida anterior en el commit \texttt{35b7a28}, según \texttt{tables\_v5.md}. Elaboración propia.
\end{table}

